\documentclass{article}
\usepackage[margin=1.5cm]{geometry}
\usepackage{amssymb}
\usepackage{float}
\usepackage{caption}

\begin{document}
\title{CSCD 327 Homework 3}
\author{Dustin Randall}
\maketitle

\section{Which of the following FDs may hold?}
\subsection{$Z \rightarrow Y$}
$Z \rightarrow Y$ cannot hold because Y at z3 = y2 and y4
\subsection{$\{W, X \} \rightarrow Y$}
$\{W, X \} \rightarrow Y$ may hold. There are no two tuples with the same $\{W, X\}$ values to compare.
\subsection{$\{Y, Z \} \rightarrow X$}
$\{Y, Z \} \rightarrow X$ may hold. There are no two tuples with the same $\{Y, Z\}$ values to compare.

\section{Is the relation in 2NF?}
No, the relation is not in 2NF because the nonprime attributes Name and $\{Course Title, Instructor\}$
are not fully dependent on the primary key $\{StudentID, CourseID\}$.

\section{CarSale(\underline{VIN}, DateSold, \underline{Salesperson}, Commission, Discount)}
\subsection{Is CarSale in 2NF?}
No, CarSale is not in 2NF because the nonprime attribute Commission is partially dependent on the primary key $\{VIN, Salesperson\}$.
\subsection{Normalize to 2NF then 3NF.}
\subsubsection{2NF}
To convert into 2NF, we need to remove the partial dependency of Commission on Salesperson.
We extract CarSale2 with Salesperson as the primary key, and Commission as the nonprime attribute. \\
CarSale1(\underline{VIN}, DateSold, \underline{Salesperson}, Discount) \\
CarSale2(\underline{Salesperson}, Commission) \\
\subsubsection{3NF}
To convert into 3NF, we need to remove the transitive dependency of Discount on VIN by creating a separate relation. \\
CarSale1(\underline{VIN}, DateSold, \underline{Salesperson}) \\
CarSale2(\underline{Salesperson}, Commission) \\
CarSale3(\underline{DateSold}, Discount) \\
\subsection{Justify each table as 3NF.}
CarSale1 is in 3NF because DateSold is fully functionally dependent on the primary key $\{VIN, Salesperson\}$ \\
CarSale2 is in 3NF because Commission is fully functionally dependent on the primary key Salesperson. \\
CarSale3 is in 3NF because Discount is fully functionally dependent on the primary key DateSold.
\subsection{Use the Chase Test to check if the decomposition is lossless.}
\begin{table}[H]{Initial Table}
\begin{tabular}{|c|c|c|c|c|c|}
    \hline
    & VIN & DateSold & Salesperson & Commission & Discount \\
    \hline
    CarSale1 & b11 & b12 & b13 & b14 & b15 \\
    \hline
    CarSale2 & b21 & b22 & b23 & b24 & b25 \\
    \hline
    CarSale3 & b31 & b32 & b33 & b34 & b35 \\
    \hline
    \end{tabular}
\end{table}
\begin{table}[H]{Apply Columns for CarSale 1-3}
\begin{tabular}{|c|c|c|c|c|c|}
    \hline
    & VIN & DateSold & Salesperson & Commission & Discount \\
    \hline
    CarSale1 & a1 & a2 & a3 & b14 & b15 \\
    \hline
    CarSale2 & b21 & b22 & a3 & a4 & b25 \\
    \hline
    CarSale3 & b31 & a2 & b33 & b34 & a5 \\
    \hline
    \end{tabular}
\end{table}

\begin{table}[H]{Apply $Salesperson \rightarrow Commission$}
\begin{tabular}{|c|c|c|c|c|c|}
    \hline
    & VIN & DateSold & Salesperson & Commission & Discount \\
    \hline
    CarSale1 & a1 & a2 & a3 & a4 & b15 \\
    \hline
    CarSale2 & b21 & b22 & a3 & a4 & b25 \\
    \hline
    CarSale3 & b31 & a2 & b33 & b34 & a5 \\
    \hline
    \end{tabular}
\end{table}

\begin{table}[H]{Apply $DateSold \rightarrow Discount$}
\begin{tabular}{|c|c|c|c|c|c|}
    \hline
    & VIN & DateSold & Salesperson & Commission & Discount \\
    \hline
    CarSale1 & a1 & a2 & a3 & a4 & a5 \\
    \hline
    CarSale2 & b21 & b22 & a3 & a4 & b25 \\
    \hline
    CarSale3 & b31 & a2 & b33 & b34 & a5 \\
    \hline
    \end{tabular}
\end{table}
Since the first row now has all a's, the decomposition is lossless.

\section {$Product(Supplier, Brand, Type, Quantity, Store)$}
\subsection{Is Product in 3NF? Why or why not?}
Product can be considered in 3NF given the 2 FDs: $\{Brand, Type\} \rightarrow Supplier$ and $\{Supplier, Store, Type\} \rightarrow Brand$.
Supplier and Brand are both prime attributes of Product.
\subsection{Is Product in BCNF? Why or why not?}
Product is not in BCNF because both FDs have non-superkey determinants. 
$\{Brand, Type\}$ and $\{Supplier, Store, Type\}$ are not superkeys of Product.

\section{$R(A, B, C, D)$ with FDs: $A \rightarrow B$, $B \rightarrow C$}
\begin{table}[H]{Initial Table}
\begin{tabular}{|c|c|c|c|c|}
    \hline
    & A & B & C & D \\
    \hline
    $\{A, B\}$ & b11 & b12 & b13 & b14 \\
    \hline
    $\{B, C\}$ & b21 & b22 & b23 & b24 \\
    \hline
    $\{C, D\}$ & b31 & b32 & b33 & b34 \\
    \hline
    \end{tabular}
\end{table}

\begin{table}[H]{Substitute As for columns}
\begin{tabular}{|c|c|c|c|c|}
    \hline
    & A & B & C & D \\
    \hline
    $\{A, B\}$ & a1 & a2 & b13 & b14 \\
    \hline
    $\{B, C\}$ & b21 & a2 & a3 & b24 \\
    \hline
    $\{C, D\}$ & b31 & b32 & a3 & a4 \\
    \hline
    \end{tabular}
\end{table}

\begin{table}[H]{Apply $A \rightarrow B$}
\begin{tabular}{|c|c|c|c|c|}
    \hline
    & A & B & C & D \\
    \hline
    $\{A, B\}$ & a1 & a2 & b13 & b14 \\
    \hline
    $\{B, C\}$ & a1 & a2 & a3 & b24 \\
    \hline
    $\{C, D\}$ & b31 & b32 & a3 & a4 \\
    \hline
    \end{tabular}
\end{table}

\begin{table}[H]{Apply $B \rightarrow C$}
\begin{tabular}{|c|c|c|c|c|}
    \hline
    & A & B & C & D \\
    \hline
    $\{A, B\}$ & a1 & a2 & a3 & b14 \\
    \hline
    $\{B, C\}$ & a1 & a2 & a3 & b24 \\
    \hline
    $\{C, D\}$ & b31 & b32 & a3 & a4 \\
    \hline
    \end{tabular}
\end{table}
At this point we've run out of FDs to apply, and no rows contain all a's, so the decomposition is lossy.

\end{document}